\section{Experiment 4: what allocation responds to}\label{sec:qualifier}

Experiment~1 contrasts a drifted memory against a clean one carrying its caveat
and reads the difference (\S\ref{sec:h2}) as detection of the omission. That
contrast varies four things at once. Appending the caveat takes the target body
to \BodyCleanTarget{} characters against \BodyMin--\BodyMax{} for every other
memory --- \BodyRatio$\times$ --- and makes it the only memory in the set
carrying a \emph{However} clause. Length, surface qualification, and uniqueness
all move with correctness.

Experiment~4 separates them. Four arms run concurrently, differing only in what
is appended to the corrupted memory's body, matched to the true caveat on
character length and clause structure: \textbf{drifted} (nothing appended,
\BodyDriftTarget{} characters), \textbf{clean-positive} (a positive elaboration,
no hedge), \textbf{clean-neutral} (a measurement-hygiene note that weakens
confidence in the finding without saying the action is risky), and
\textbf{clean-negative} (the true caveat). Four arms $\times$ six models
$\times$ 25 gives \FourAllN{} episodes. Only the first pass is called, since
first-pass verification is the pre-registered outcome and the post-retrieval
turn bears on none of H12--H13.

\paragraph{Intent is a mediator here, not a confound.} The arms differ in
first-pass intent: the positive elaboration makes pricing more attractive
(\FourPosIntent/\FourPosAllN{} episodes intend it, against
\FourDriftIntent/\FourDriftAllN{} under drift) and the true caveat removes the
plan entirely (\FourNegIntent/\FourNegAllN). Because the qualifier is
randomised and intent is measured after it, that difference is an effect of the
treatment, not a pre-treatment imbalance. Conditioning on it would be
post-treatment selection and could bias the estimate rather than clean it.

We therefore keep the pre-registered marginal contrasts as the primary result;
they are the randomised total effect of qualifier content on verification. All
three are supported: clean-neutral \FourNeutNegPreA{} against clean-negative
\FourNeutNegPreB{} ($\FourNeutNegPreCMH$), clean-positive \FourPosNegPreA{}
($\FourPosNegPreCMH$), and clean-neutral against drifted \FourNeutDriftPreA{}
versus \FourNeutDriftPreB{} ($\FourNeutDriftPreCMH$).

We also report the intent-misaligned subset, because it shows the allocation
pattern among episodes that share the same observed intent, which is what makes
the mechanism legible. It is descriptive: selected after observing the data,
conditioned on a post-treatment variable, therefore open to selection bias, and
not a less confounded estimate of anything.
Both sets are in Table~\ref{tab:qualifier}, and every contrast has the same
sign in both.

\begin{table}[t]
\centering\small
\caption{Experiment 4. Verification of the corrupted memory, first pass, budget
two, by qualifier content. Bodies are matched on length and clause structure
across the three appended arms; only the content differs. The top row is the
randomised contrast over all \FourAllN{} episodes and is the primary result.
The bottom row is the subset of episodes in which the agent does not intend the
action that memory backs --- a post-treatment condition, so descriptive only,
not a randomised conditional effect.}
\label{tab:qualifier}
\begin{tabular}{lcccc}
\toprule
& clean-negative & drifted & clean-neutral & clean-positive \\
& (true caveat) & (no qualifier) & (irrelevant hedge) & (positive) \\
\midrule
\emph{randomised (primary)} & \FourNegAllK/\FourNegAllN{} & \FourDriftAllK/\FourDriftAllN{} & \FourNeutAllK/\FourNeutAllN{} & \FourPosAllK/\FourPosAllN{} \\
 & (\FourNegAllPct\%) & (\FourDriftAllPct\%) & (\FourNeutAllPct\%) & (\FourPosAllPct\%) \\
\addlinespace
\emph{intent-misaligned subset (descriptive)} & \FourNegK/\FourNegN{} & \FourDriftK/\FourDriftN{} & \FourNeutK/\FourNeutN{} & \FourPosK/\FourPosN{} \\
 & (\FourNegPct\%) & (\FourDriftPct\%) & (\FourNeutPct\%) & (\FourPosPct\%) \\
\bottomrule
\end{tabular}
\end{table}

Table~\ref{tab:qualifier} orders the arms in a way a salience account does not
predict. The true caveat does not raise verification; it \emph{collapses} it to
\FourNegPct\%. That figure is pooled, and the suppression is not universal:
it holds in \FourSuppressModels{} of six models, while in Sonnet
(\FourSonnetNeg{} with the caveat against \FourSonnetDrift{} without) and Haiku
(\FourHaikuNeg{} against \FourHaikuDrift) a stated caveat draws slightly more
attention rather than less. A written constraint is not universally a reason to
stop looking. A matched-length qualifier that raises a question without
answering it raises verification (\FourNeutPct\% and \FourPosPct\%). The
silently drifted memory sits between them at \FourDriftPct\%.

The mechanism is visible in where the budget goes instead. In the arm carrying
the true caveat, one model's first-pass intent was onboarding in every episode
and its two credits went to the memories backing that plan; the corrupted memory
was not requested once in 25 episodes. Nothing about it had become less
important --- its constraint was already written down, so there was nothing to
recover.

\paragraph{Experiment 1's H2 result holds; its interpretation does not.}
Verification really is higher when the caveat is absent than when it is present,
and that survives intent-matching (\FourDriftK/\FourDriftN{} against
\FourNegK/\FourNegN{}, $\FourDriftNegCMH$). But it is not the agent noticing
that something was removed, and \S\ref{sec:h2} should not have read it that way.
Detection of the omission is not required to produce the effect, and Experiment~4
provides no evidence that it occurs.

\paragraph{Silence is stealthier than qualification.} A memory that hedges about
something irrelevant draws more scrutiny than one whose material caveat was
deleted (\FourNeutDriftDiff{} points, $\FourNeutDriftCMH$). Consolidation that
drops a clause does not merely remove information --- it removes the surface cue
that would have made the record look worth checking. The record becomes
cleaner-looking and correspondingly less audited, which is the worst combination
for a store that later sessions inherit.

\section{Experiment 5: does benign consolidation produce this?}\label{sec:drift}

Experiments 1--4 install the corruption by hand, which buys clean ground truth.
It also assumes what \S1 assumed: that summarization drops caveats. We tested
that assumption on one consolidation procedure and did not get confirmation.
What follows describes the chain below, not summarization systems generally.

We ran a consolidation chain over the same six source records with no
instruction to preserve or drop anything and no mention of caveats, qualifiers,
or what matters: an analyst writes session notes; the notes are consolidated
into durable lessons; durable memory is then re-summarized four times as it ages
and later cycles compete for context, with targets stepping
$130 \rightarrow 110 \rightarrow 90 \rightarrow 70$ characters. \FiveChains{}
chains, \FiveGenerations{} generations each, six models,
\FiveScored{} scored (chain, memory) pairs.

Survival is judged by a model that sees the original source record and all
generations at once and reports, for each, whether the source's specific
constraint is still stated. Two of the six records carry no qualifier and act as
a false-positive check: the scorer identified them as unqualified in every case
and reported a constraint in \FiveFalsePosAll{} judgments where none exists.

\paragraph{The premise fails for negative outcomes.} At the tightest
compression the retention harm is still stated in \FiveNegResDPct\% of chains,
ranging \FiveNegPerModelLo{} to \FiveNegPerModelHi{} across models. Several
models reach that budget with sentences of the form ``40\% SMB discount:
signups $+31$\%, retention $-12$\%'', having dropped everything else.

\paragraph{What erodes is scope, and normative force.} Scope restrictions go
roughly twice as fast: \FiveScopeResDPct\% survive at the same compression
against \FiveNegResDPct\%. Within the corrupted memory itself, separating the
finding from the instruction shows the split directly --- at consolidation
\FiveNumCons{} chains state the negative outcome and \FiveProhCons{} state a
prohibition; after four re-summarizations \FiveNumResD{} still state the outcome
and \FiveProhResD{} state a prohibition. What consolidation removes is not the
fact that discounting hurt retention. It is the sentence saying not to do it
anyway. (This split uses keyword lists written after reading the generations and
is exploratory; the survival scoring above is the pre-registered measure.)

\paragraph{A pre-registered hypothesis we could not run, which is the result.}
H15 asked whether naturally drifted lessons --- chains where the caveat is gone
and the surviving sentence is still true --- reproduce the intent-aligned
allocation result. It is not runnable. A body matching the hand-constructed
drifted one, carrying no negative content at all, appears in
\FivePurelyPositive{} chain-generations. Not once. Every generation of every
chain keeps some negative signal about the discount, even at seventy
characters. We report H15 as unrunnable rather than substituting a softer test
for it, because the reason it cannot be run is the finding.

\subsection{Exploratory: allocation under a realistically drifted memory}\label{sec:natalloc}

What consolidation does produce is a lesson that keeps the number and loses the
prohibition. We ran the Experiment-1 allocation protocol on those bodies
(\NatChains{} distinct chains, \NatLenMin--\NatLenMax{} characters, sampled by
seed and recorded per episode; \NatN{} episodes). This was written after the
Experiment-5 results were known and is exploratory throughout.

It points the opposite way from the hand-constructed drift. The naturally
drifted memory is verified in \NatK/\NatN{} episodes (\NatPct\%) against
\FourDriftAllK/\FourDriftAllN{} (\FourDriftAllPct\%) for the hand-constructed
one under the identical protocol, and \emph{no} episode --- \NatPricingIntent{}
of \NatN{} --- intends promotional pricing at all.

That is Experiment~4's mechanism appearing again. \NatPct\% sits beside the
true caveat's \FourNegAllPct\%, not beside the hand-drifted
\FourDriftAllPct\%: a memory that still says ``retention $-12$\%'' is treated
as a memory whose constraint is stated, because it is. The retained number
carries enough of the caveat to keep the agent off the action, so there is
nothing for the budget to recover.

\paragraph{What this costs us.} The vulnerability Experiments 1--4 measure
requires a corruption that removes the negative content entirely. Our
consolidation chain never produced one. So the allocation signature is real and
its downstream behavioural consequence is causally established; on the evidence
here, however, the particular corruption that exposes it has to arise somewhere
other than the consolidation chain we tested --- an aggressive rewrite, a lossy
migration, an adversary, or a decision geometry where the surviving number does
not settle the question. Claiming that ordinary
summarisation walks agents into this failure would go beyond what we measured,
and \S1's original framing did.

\paragraph{Factual preservation and constraint loss are different failures.}
The distinction this experiment forces is between a stored claim becoming
\emph{false} and a stored claim staying true while the conditions under which it
should guide action disappear. Only the second occurs here, and it is the harder
one to notice: nothing in the record is wrong, the numbers are intact and
checkable, and what is missing is the scope the finding applied to and the
instruction not to generalise it. A memory can stay factually accurate and become
decision-dangerous --- and \S\ref{sec:qualifier} says such a record draws
\emph{less} verification, because a lesson that keeps its numbers looks
well-evidenced. A memory system that validates stored facts would pass it.

\paragraph{Consequences for our own threat model.} Deleting an entire negative
caveat is not what the consolidation process we tested produces. A benign-compression threat model
should be built on the operators this compression actually applied: loss of scope,
and loss of normative constraint while the supporting number survives. Those are
more dangerous than they sound, because a lesson that retains its numbers looks
well-evidenced --- and Experiment~4 says such a record draws less verification,
not more.

The allocation results are unaffected: they concern which beliefs an agent
checks given a corrupted one, and hold for any corruption with the same decision
geometry. What does not hold is the claim that this particular corruption arises
on its own.

\paragraph{A pre-registered rule that failed.} We registered that two scorers, a
keyword rule and a model scorer, must agree. The rule is mis-specified and we
report that rather than its output: disagreements occur precisely where one
scorer sees a lost qualifier, so conditioning on agreement censors informatively
and returns 100\% survival at every generation by construction. The keyword rule
also has very low specificity on scope restrictions, firing on any mention of
the segment term, so it bounds rather than estimates. We report the model
scorer, validated against the unqualified controls, as primary.
